\chapter{LITERATURE REVIEW}

\section{Introduction to the Literature Review}

The literature surrounding AI-driven recruitment, retrieval-augmented generation (RAG), vector representations, code language models, and multi-agent orchestration spans multiple distinct domains within computer science, natural language processing (NLP), and empirical software engineering \cite{ref1, ref4, ref33, ref44}. As artificial intelligence becomes deeply embedded in enterprise recruitment workflows, understanding the historical progression, theoretical foundations, and operational vulnerabilities of candidate evaluation systems is essential for developing robust credential authentication frameworks \cite{ref2, ref3, ref23}.

This chapter synthesizes foundational literature across six key research areas:
\begin{enumerate}[leftmargin=*, label=\arabic*.]
    \item \textbf{AI-Driven Resume Screening Systems:} Examining the evolution from deterministic keyword matching and classic machine learning classifiers (SVMs, Naive Bayes) to transformer-based entity extraction and scoring \cite{ref1, ref9, ref11, ref7, ref8}.
    \item \textbf{Retrieval-Augmented Generation (RAG):} Analyzing dense knowledge retrieval architectures (DPR, FiD, RETRO, Self-RAG, Active RAG) and their role in mitigating large language model hallucination \cite{ref2, ref5, ref30, ref33, ref45, ref46}.
    \item \textbf{Sentence Embeddings and Vector Databases:} Reviewing dense vector space mapping (Sentence-BERT, Instructor, MTEB, AnglE) and sub-millisecond nearest-neighbor search engines (ChromaDB, Pinecone, HNSW graphs) \cite{ref31, ref32, ref47, ref34, ref41, ref48, ref49}.
    \item \textbf{Large Language Models and Code Analysis:} Surveying general foundation LLMs (GPT-3, GPT-4, LLaMA, Mistral, PaLM, Gemini) and specialized code LLMs (Codex, StarCoder, Code Llama, CodeBERT, UniXcoder, CodeT5+) alongside Chain-of-Thought (CoT) prompting techniques \cite{ref9, ref20, ref30, ref16, ref37, ref50, ref51, ref52, ref53, ref17, ref54, ref55, ref56, ref57, ref58, ref42, ref59, ref60}.
    \item \textbf{Multi-Agent LLM Frameworks:} Investigating autonomous agent coordination paradigms (Generative Agents, AutoGen, CAMEL, MetaGPT, ChatDev) and Standard Operating Procedures (SOPs) \cite{ref61, ref62, ref38, ref63, ref43, ref64, ref39}.
    \item \textbf{Empirical Software Engineering \& Ethical AI Hiring:} Evaluating GitHub repository mining standards, developer contribution metrics, and regulatory compliance frameworks (EU AI Act, UNESCO Recommendations) \cite{ref6, ref15, ref25, ref24, ref14, ref26, ref22, ref36, ref65, ref66}.
\end{enumerate}

Finally, research gaps in existing literature are identified, establishing the theoretical and practical foundation for CodeAudit AI \cite{ref20, ref23, ref8}.

\section{AI-Driven Resume Screening Systems}

Early automated resume screening systems emerged in the late 1990s and early 2000s, driven by the need to manage massive applicant volumes generated by online job portals \cite{ref1, ref7}. These foundational Applicant Tracking Systems (ATS) relied on deterministic regular expressions and keyword frequency counting to rank candidate resumes against job description keywords \cite{ref50, ref8}. Machine learning approaches subsequently introduced supervised learning models—such as Support Vector Machines (SVMs), Naive Bayes classifiers, and Decision Trees—trained on Term Frequency-Inverse Document Frequency (TF-IDF) feature vectors to categorize candidates into predefined suitability brackets \cite{ref7, ref13}.

However, keyword-centric machine learning models suffered from severe operational flaws \cite{ref9}. Because TF-IDF and bag-of-words representations treat words as independent tokens without semantic context, these systems were easily exploited through keyword stuffing \cite{ref10, ref19}. Candidates who artificially inserted technical buzzwords into invisible text blocks or overloaded skills sections scored highly, while candidates who described equivalent project experience using different vocabulary were misclassified and discarded \cite{ref11, ref12}. Furthermore, classic classifiers possessed zero capacity to verify whether listed skills reflected actual practical capability or mere textual familiarity \cite{ref2, ref3}.

The introduction of transformer architectures, specifically BERT (Bidirectional Encoder Representations from Transformers) \cite{ref12}, fundamentally shifted NLP approaches in recruitment \cite{ref11}. Researchers fine-tuned BERT models for Named Entity Recognition (NER) to extract technical skills, job titles, and educational institutions from unformatted PDF text streams \cite{ref9, ref19}. Subsequent work utilized sentence transformer embeddings to compute Cosine Similarity scores between parsed resume vectors and job requirement embeddings \cite{ref41, ref8}.

Despite these algorithmic advances, modern transformer-based resume screeners evaluate resumes strictly as self-contained textual claims \cite{ref23, ref19}. As highlighted by Sinha et al. \cite{ref8} and Faliagka et al. \cite{ref7}, evaluating candidate claims in isolation without cross-referencing external implementation evidence leaves recruitment platforms vulnerable to candidate skill exaggeration, fabricated experience, and AI-generated resume stuffing \cite{ref4, ref5, ref17}.

Recent work by Lo et al. \cite{ref23} explored multi-agent LLM systems to evaluate candidate resumes across multiple qualitative dimensions, such as experience depth and leadership potential. However, their system evaluates resumes solely against job description prompts rather than cross-referencing skill claims against functional source code repositories \cite{ref20, ref23}. Consequently, a fundamental verification gap persists between claimed technical qualifications and actual coding capability \cite{ref15, ref16, ref24}.

\section{Retrieval-Augmented Generation (RAG)}

Retrieval-Augmented Generation (RAG) was formally introduced by Lewis et al. \cite{ref32} to address the inherent parametric memory limitations of Large Language Models. Foundation LLMs store knowledge implicitly within their neural weights, making them susceptible to factual hallucinations, knowledge cut-off boundaries, and an inability to inspect private or domain-specific document collections \cite{ref9, ref17, ref18}. RAG resolves these challenges by coupling an autoregressive language model with an external, non-parametric vector index, retrieving relevant document passages at inference time to ground model outputs in empirical evidence \cite{ref33, ref45, ref32}.

Subsequent research introduced architectural enhancements to traditional RAG pipelines \cite{ref33}. Izacard and Grave \cite{ref45} developed Fusion-in-Decoder (FiD), which processes retrieved passages independently through transformer encoder layers before concatenating their representations in the decoder, enabling scalable multi-document reasoning. Borgeaud et al. \cite{ref5} presented RETRO (Retrieval-Enhanced Transformer), demonstrating that retrieving over trillions of tokens allows compact language models to achieve performance comparable to parameter-heavy frontier models.

Gao et al. \cite{ref33} presented a comprehensive RAG survey, categorizing systems into Naive RAG, Advanced RAG (incorporating pre-retrieval query rewriting and post-retrieval reranking), and Modular RAG (enabling dynamic routing and multi-agent iteration). Self-RAG \cite{ref2} introduced self-reflection tokens that allow language models to dynamically decide when to retrieve passages and evaluate output factual correctness. Active RAG architectures, such as FLARE \cite{ref46} and REPLUG \cite{ref67}, dynamically interleave document retrieval with text generation based on token-level confidence thresholds \cite{ref68, ref69}.

In technical resume verification, RAG provides the mathematical foundation for anchoring LLM judgment in concrete codebase evidence \cite{ref31, ref33, ref32}. Rather than prompting an LLM to guess whether a candidate knows a claimed skill, CodeAudit AI utilizes RAG to retrieve the top-5 most semantically relevant source code chunks from ChromaDB, forcing the LLM Judge Agent to evaluate actual code syntax before issuing a verdict \cite{ref10, ref27, ref25, ref29}.

\section{Sentence Embeddings and Vector Databases}

Transforming natural language queries and source code snippets into dense, semantically rich vector representations is essential for effective similarity search \cite{ref70, ref41, ref49}. Reimers and Gurevych \cite{ref41} introduced Sentence-BERT (SBERT), utilizing Siamese and triplet network structures to derive semantically meaningful sentence embeddings that can be efficiently compared using Cosine Similarity or Euclidean distance metrics.

Subsequent research focused on contrastive learning and instruction-tuned embedding models \cite{ref71, ref48, ref49}. Su et al. \cite{ref48} introduced Instructor embeddings, demonstrating that prefixing domain-specific task instructions to text inputs significantly improves cross-domain retrieval performance. Wang et al. \cite{ref49} presented contrastive learning frameworks that established state-of-the-art benchmarks on the Massive Text Embedding Benchmark (MTEB) \cite{ref47}. Li et al. \cite{ref71} introduced AnglE embeddings, optimizing gradient loss metrics to capture fine-grained semantic nuances in text and code structures. In CodeAudit AI, the \texttt{all-MiniLM-L6-v2} sentence transformer \cite{ref41} is selected for its balance of high semantic fidelity (384 dimensions) and low computational latency.

Vector database architectures have evolved rapidly to support sub-millisecond similarity search over millions of high-dimensional vectors \cite{ref31, ref34}. Modern vector stores—such as ChromaDB \cite{ref31} and Pinecone \cite{ref34}—utilize Hierarchical Navigable Small World (HNSW) graph algorithms and Inverted File (IVF) indexing to perform approximate nearest-neighbor (ANN) search. Crucially, ChromaDB supports metadata filtering, allowing vector similarity queries to be constrained by metadata attributes \cite{ref31, ref33}. In CodeAudit AI, this metadata filtering mechanism is leveraged to enforce strict candidate-level data isolation during vector retrieval \cite{ref33, ref34}.

\section{Large Language Models and Code Analysis}

Foundation language models have demonstrated remarkable reasoning capabilities across diverse tasks \cite{ref9, ref30, ref16, ref17, ref56, ref18}. Autoregressive transformer models scaled to tens or hundreds of billions of parameters exhibit emergent capabilities, including zero-shot instruction following, complex reasoning, and structured output generation \cite{ref9, ref52, ref42}. Frontier models such as GPT-3 \cite{ref9}, GPT-4 \cite{ref17}, LLaMA 3 \cite{ref16}, Mistral 7B \cite{ref51}, PaLM \cite{ref30}, and Gemini \cite{ref55} serve as the core intelligence engines for autonomous agent frameworks \cite{ref38, ref64, ref44}.

In software engineering, specialized code LLMs have been trained on vast open-source code repositories \cite{ref3, ref20, ref53, ref54, ref58}. Chen et al. \cite{ref20} introduced Codex, demonstrating automated code generation and natural language to code translation. Li et al. \cite{ref53} presented StarCoder, trained on over 80 programming languages from GitHub. Roziere et al. \cite{ref54} developed Code Llama, providing enhanced code completion, infilling, and natural language explanation capabilities. Parallel research in code representation learning produced models such as CodeBERT \cite{ref37}, UniXcoder \cite{ref50}, and CodeT5+ \cite{ref58}, establishing structural representations of Abstract Syntax Trees (ASTs) and data-flow graphs \cite{ref3, ref37, ref50}.

Chain-of-Thought (CoT) prompting \cite{ref42} further enhanced LLM reasoning by encouraging models to generate intermediate step-by-step reasoning chains before producing a final answer. Kojima et al. \cite{ref52} demonstrated zero-shot CoT using simple prompt prefixes. Wang et al. \cite{ref57} introduced Self-Consistency, sampling multiple CoT reasoning paths and selecting the majority verdict. Yao et al. \cite{ref59} presented Tree of Thoughts (ToT), enabling tree-based search over intermediate reasoning steps, while Zhou et al. \cite{ref60} explored Least-to-Most prompting for complex task decomposition. In CodeAudit AI, CoT prompting guides the LLM Judge Agent to methodically verify retrieved code snippets against claimed technical skills \cite{ref10, ref25, ref42, ref29}.

\section{Multi-Agent LLM Frameworks}

Multi-agent coordination frameworks enable specialized autonomous agents to collaborate on complex problem-solving tasks that exceed the capacity of a single monolithic prompt \cite{ref61, ref62, ref38, ref63, ref43, ref39}. Park et al. \cite{ref38} demonstrated generative agents capable of believable social coordination. Talebirad and Nadiri \cite{ref43} synthesized core principles of multi-agent collaboration, emphasizing role specialization, structured message passing, and environment feedback loops \cite{ref38, ref39}.

Open-source multi-agent frameworks have formalized these design patterns:
\begin{itemize}[leftmargin=*]
    \item \textbf{AutoGen \cite{ref39}:} Enables conversational multi-agent workflows with customizable human-in-the-loop controls.
    \item \textbf{CAMEL \cite{ref62}:} Formulates communicative agent setups using role-playing prompts to solve complex tasks autonomously.
    \item \textbf{MetaGPT \cite{ref61}:} Embeds software engineering Standard Operating Procedures (SOPs) into agent prompt sequences, assigning distinct roles (Product Manager, Architect, Engineer) to produce structured software artifacts.
    \item \textbf{ChatDev \cite{ref63}:} Models software development as an iterative chat-chain between virtual personas collaborating across design, coding, testing, and documentation phases.
\end{itemize}

Comprehensive surveys by Wang et al. \cite{ref64} and Xi et al. \cite{ref44} confirm that multi-agent architectures outperform monolithic LLMs by isolating functional responsibilities, reducing context window clutter, and preventing single-point reasoning failures across processing phases \cite{ref20, ref23, ref64, ref44}.

\section{GitHub as an Empirical Data Source and Ethical AI Hiring}

Mining open-source software repositories on GitHub provides invaluable empirical insight into practical software engineering skills \cite{ref6, ref15, ref25, ref24, ref26, ref36}. Kalliamvakou et al. \cite{ref24} established methodology guidelines for mining GitHub data, noting that public commits, repository structures, language distributions, and dependency manifests reflect real-world implementation experience \cite{ref6, ref25, ref26}. Vasilescu et al. \cite{ref36} and Borges et al. \cite{ref6} analyzed developer contribution patterns and project popularity metrics, confirming that public code artifacts serve as authentic skill indicators.

Concurrently, ethical considerations and regulatory compliance in automated recruitment have drawn intense academic and legal attention \cite{ref4, ref14, ref22, ref65, ref66}. Raghavan et al. \cite{ref22} evaluated bias mitigation strategies in hiring algorithms, warning against proxy discrimination in automated resume screeners. Mujtaba and Mahapatra \cite{ref14} emphasized the need for auditable, transparent decision-making in HR technology.

International regulatory frameworks—such as the EU Artificial Intelligence Act \cite{ref65} and UNESCO Recommendations on the Ethics of AI \cite{ref66}—mandate strict transparency, auditability, data privacy, and human oversight in employment-related AI systems. CodeAudit AI directly addresses these compliance mandates by replacing subjective scoring with deterministic, evidence-grounded verification backed by explicit source code citations \cite{ref2, ref65, ref66}.

\section{Limitations and Challenges in Previous Research Work}

A critical examination of the literature reveals several fundamental methodological, architectural, and operational limitations in existing recruitment and code intelligence frameworks:

\begin{enumerate}[leftmargin=*, label=\arabic*.]
    \item \textbf{Vulnerability to Keyword Inflation and Adversarial Stuffing:} Foundational Applicant Tracking Systems (ATS) and classic TF-IDF/SVM models operate under the naive assumption that keyword occurrences correlate directly with technical competence \cite{ref1, ref7, ref11}. Consequently, these systems are easily manipulated through keyword stuffing and prompt injection, allowing unqualified candidates to game recruitment filters while discarding qualified practitioners who use alternate terminology \cite{ref10, ref19}.
    
    \item \textbf{Absence of Empirical Code Grounding:} State-of-the-art transformer-based resume screeners \cite{ref12, ref23} evaluate candidate qualifications purely as self-contained textual claims. As LLMs have democratized the generation of polished, AI-assisted resumes, text-only screening models cannot distinguish authentic software engineering experience from synthetic or exaggerated claims \cite{ref4, ref8, ref17}.
    
    \item \textbf{Semantic Inability of Static Code Analyzers:} While traditional static analysis tools (e.g., SonarQube, Bandit) assess code syntax, linting errors, and security vulnerabilities \cite{ref24, ref36}, they lack natural language understanding and contextual reasoning capabilities. They cannot map unstructured resume claims (e.g., ``Architected scalable microservices'') to specific functional abstractions within a codebase \cite{ref15, ref20}.
    
    \item \textbf{Boilerplate Scaffolding Blindspots in Dense Code Retrieval:} Generic dense code search models (e.g., CodeBERT \cite{ref37}, UniXcoder \cite{ref50}) compute vector similarities based on token co-occurrence and semantic representations. However, empirical experiments reveal that dense retrieval models frequently misclassify untouched boilerplate templates (such as default \texttt{create-react-app} scaffolds or cloned repositories) as verified implementations because they lack Abstract Syntax Tree (AST) functional filtering and Git commit authorship attribution \cite{ref25, ref26, ref58}.
    
    \item \textbf{Multi-Tenant Vector Isolation and Hallucination Risks:} Monolithic RAG architectures that feed unpartitioned code contexts into general-purpose LLMs suffer from high hallucination rates and cross-tenant data leakage risks in multi-candidate enterprise environments \cite{ref31, ref33, ref34}.
\end{enumerate}

Table \ref{tab:system_limitations_comparison} presents a comprehensive comparative analysis of existing research paradigms, summarizing their architectural limitations and highlighting the necessity of the proposed CodeAudit AI framework.

\begin{table}[htbp]
\setlength{\tabcolsep}{6pt}
\renewcommand{\arraystretch}{1.3}
\centering
\small
\caption{Comprehensive Comparison of Limitations and Challenges in Previous Research vs. CodeAudit AI}
\label{tab:system_limitations_comparison}
\begin{tabularx}{\textwidth}{@{}p{3.0cm}p{2.8cm}p{3.6cm}p{4.2cm}@{}}
\toprule
\textbf{Research Paradigm} & \textbf{Representative Systems} & \textbf{Key Capabilities} & \textbf{Critical Limitations \& Gaps} \\ \midrule
Keyword \& Classic ML ATS & Taleo, Workday ATS, Alamri et al. \cite{ref1} & Keyword frequency, TF-IDF ranking & Blind to semantic context; easily fooled by keyword stuffing; zero code verification. \\ \midrule
Static Code Analyzers & SonarQube, ESLint, Bandit \cite{ref24, ref36} & Syntax validation, linting, security scanning & Lacks NLP reasoning; cannot cross-reference or verify natural language resume claims. \\ \midrule
LLM Resume Summarizers & Deshpande et al. \cite{ref11}, Lo et al. \cite{ref23} & Entity extraction, qualitative candidate scoring & Evaluates text in isolation; vulnerable to AI-generated resume inflation and hallucination. \\ \midrule
Dense Code Retrieval & CodeBERT \cite{ref37}, UniXcoder \cite{ref50} & Semantic code vector search & High recall but blind to boilerplate scaffolds and forked repos without authorship checks. \\ \midrule
\textbf{CodeAudit AI (Proposed Work)} & \textbf{Multi-Agent AST RAG (This Research)} & \textbf{AST-guided chunking, Git authorship checks, Multi-Agent cross-examination} & \textbf{Addresses all gaps: Evidence-grounded, anti-boilerplate, 100\% multi-tenant isolated.} \\ \bottomrule
\end{tabularx}
\end{table}

\section{Research Gaps and Summary}

The literature review and comparative analysis demonstrate three paramount research gaps:
\begin{enumerate}[leftmargin=*, label=\arabic*.]
    \item \textbf{Verification Gap:} Existing screening platforms evaluate resume text without verifying claims against functional source code evidence \cite{ref1, ref9, ref8}.
    \item \textbf{Isolation Gap:} Vector retrieval systems in hiring platforms lack strict candidate-level metadata isolation, creating risks of cross-tenant data leakage \cite{ref31, ref33, ref34}.
    \item \textbf{Hallucination Gap:} Monolithic scoring models generate subjective scores without citing verifiable codebase evidence \cite{ref20, ref17, ref29}.
\end{enumerate}

By combining multi-agent orchestration, AST-guided RAG indexing, Git authorship attribution, and candidate-isolated ChromaDB storage, CodeAudit AI bridges these research gaps, establishing an empirical framework for technical skill authentication \cite{ref30, ref37, ref33, ref38}.